\documentclass[12pt]{article}

\usepackage[margin=1in]{geometry}
\usepackage[T1]{fontenc}
\usepackage[utf8]{inputenc}
\usepackage{lmodern}
\usepackage{enumitem}
\usepackage{pdfpages}
\usepackage{hyperref}
\usepackage[round,authoryear]{natbib}

\hypersetup{
  hypertexnames=false,
  colorlinks=true,
  linkcolor=black,
  citecolor=black,
  urlcolor=blue
}

\setcounter{tocdepth}{3}

\title{\textbf{The Ethical Violation Issues in Product's Environmental Advertising}}
\author{Secary Ma}
\date{22 Nov 2024}

\begin{document}

\pagenumbering{roman}

\section*{Cover Letter}

Secary Ma\\
Data Analyst\\
22 Nov 2024

\vspace{1em}

John le Blanc\\
Chief Executive Officer

\vspace{1em}

Antony Cox\\
Chief Marketing Officer

\vspace{1em}

Dear Mr. John le Blanc and Mr. Antony Cox,

This is Secary Ma, a data analyst from Product Operations Department. I have realised that there might be some ethical violations in our product's advertisement. They are mainly about using environmentally friendly raw materials in our products. Since I am a user of our own product and I have been sick recently, I went to my doctor and found that the reason for my sickness is caused by our product, with a kind of harmful ingredient in it which is not environmentally friendly. I have made further investigations and now I believe that we are using false statements in our advertisement which could do harm to both consumers and the environment.

The attached file is a report about the ethical violations, including their scope and impacts, and also what actions we could take now. As an employee and a consumer, I sincerely hope the company could put an emphasis on the ethical violations and take actions on behalf of all consumers' welfare.

I hope you could reply as soon as possible for this issue.

\vspace{1em}

Yours sincerely,\\
Secary Ma

\clearpage

\begin{titlepage}
  \centering
  \vspace*{\fill}
  {\LARGE\bfseries The Ethical Violation Issues in Product's Environmental Advertising\par}
  \vspace{2em}
  {\large Secary Ma\par}
  \vspace{0.5em}
  {\large 22 Nov 2024\par}
  \vspace*{\fill}
\end{titlepage}

\tableofcontents
\clearpage
\pagenumbering{arabic}

\section{Executive Summary}

This report focuses on the ethical violations in the company's product environmental advertisement, which involves false statements about using environmentally friendly materials to manufacture products when the products do not actually use these materials. These violations may lead to risks involving reputation, finance and compliance, and they can also cause consumer trust backlash.

Such ethical breaches, often referred to as ``greenwashing'', can raise consumer dissatisfaction and harmfully affect the environment. For the further development of our company, I sincerely suggest that we should take actions to prevent these consequences from happening.

\clearpage

\section{Introduction}

We have noticed that consumers' demands for environmentally friendly products have been continually increasing in recent years. As a large corporation established with environmental protection as our core competence, we are deeply committed to integrating sustainability into every aspect of our operations and we have a solid brand image based on the authenticity and reliability of our products. However, as the investigation of our company's products deepens, I have found that part of the raw materials of our products are not meeting environmental standards. This is obviously not the same as our advertisements claiming they are 100 percent environmentally friendly. Behaviours like this not only breach the principles of marketing integrity but also may cause legal risks and consumer trust backlash.

This report will concentrate on the scope and influence of this kind of ethical violation. First, we will discuss the definition of ethical violations in this issue, including losing trust from consumers and lack of transparency in advertising. Next, we will discuss the scope and impacts that would be caused by this issue. The scope part will contain the number of influenced consumers and the effects on the environment caused by this issue. In the impact part we will mainly discuss how this case would affect our consumers and our company.

Then the report will provide some practical and workable suggestions to make damage control and improve our manufacturing processes. We firmly believe that the best products must be environmentally friendly, so we must make a commitment to the public and environment to provide sustainable, harmless products without troubled materials.

\section{Discussion}

First, we should hold a discussion about the ethical violations in this circumstance, which should contain the definition, scope and impacts of the violations in this case.

\subsection{The Definition of the Ethical Violations}

The company has violated the ethical principles of truthfulness and transparency in its environmental product advertising. They are specifically presented as:

\begin{itemize}
  \item using non-environmentally friendly materials to manufacture part of the products;
  \item using false statements of environmental performance of the product in marketing, which could mislead consumers.
\end{itemize}

This behaviour has breached the industry ethical standards and may further violate consumer protection law and advertising regulations. It has been claimed that an enterprise should adhere to honesty and transparency principles during marketing, especially when making claims about the features of its products \citep{steiner2014business}.

\subsection{The Scope and Impacts of the Ethical Violations}

\subsubsection{Scope}

\paragraph{Consumers}
\begin{itemize}
  \item We had 30 percent of total products in the last year using troubled materials which have been sold to customers nationwide.
  \item Among these troubled products, 50 percent were using completely non-recyclable materials and the other 50 percent consisted of environmentally harmful ingredients.
  \item The approximate number of customers is about 15,000 people.
\end{itemize}

\paragraph{Environment}
\begin{itemize}
  \item The product waste could pollute the marine environment and impact a large number of marine life species.
  \item The use of non-biodegradable raw materials may lead to the overexploitation of resources. For example, some raw materials may come from scarce natural resources, resulting in excessive reliance on and destruction of ecosystems such as forests and minerals, which could lead to an increase in these resources' prices \citep{hotelling2021exhaustible} and less profit in our productions.
\end{itemize}

\subsubsection{Impacts}

\paragraph{To consumers}
Our false advertisement has misled consumers into purchasing our products based on incorrect information. According to the research by \citet{torelli2019greenwashing}, we can determine that the misleading could cause consumer backlash and violations to their rights and interests.

\paragraph{To the corporation}
The impacts to the corporation could be divided into three kinds of risks, which are reputation, finance and compliance.

\begin{description}[leftmargin=2.8cm,style=nextline]
  \item[Reputation risks] As the non-environmentally friendly material use issue is being exposed, our corporation may face strong criticism from consumers and journalists.
  \item[Finance risks] Our corporation may face a heavy fine for breaching environmental protection laws, and meanwhile we may provide compensation to consumers who had been influenced by the troubled products.
  \item[Compliance risks] The false statements in advertising have violated the advertising laws and may lead to investigations by regulatory authorities.
\end{description}

Next, we still need to hold a discussion of the public's opinions of false environmental advertising issues. It is necessary to understand the shift in our position regarding public trends before and after the issue.

\subsection{Public's Opinions}

According to the industry data, global consumers' attention to sustainable products has increased in recent years, which means their tolerance of products with non-environmentally friendly materials is becoming lower and lower. Besides, false advertising, which is also known as greenwashing, is leading to less public trust \citep{greenpeace2021greenwashers}.

Therefore, we should avoid manufacturing products with non-environmentally friendly materials, take damage control actions such as communicating transparently with consumers or recalling the issued products to prevent the violations causing further harmful effects.

\section{Recommendations}

\subsection{Transparent Communication with Consumers}

As we mentioned earlier, we must take steps to prevent further harmful consequences resulting from these violations. We should acknowledge that some of the materials used in our products are not genuinely environmentally friendly, and we are sincerely concerned about the implications and potential outcomes of this issue. We should communicate all the details involved to demonstrate our sincere apologies for what we have done and to regain public confidence.

\subsection{Recall the Troubled Products as Soon as Possible}

Since we had manufactured the issued products, we should recall the products as many as we can and as soon as possible. A one-minute delay in response could result in the death of a living creature due to our issue. To be responsible for our consumers, we should recall the sold products and provide compensation to those who purchased our products but may lose trust once they learn the truth.

\subsection{Set Up Consumer's Complaint and Feedback System}

We do not know how many consumers have realised their purchased products have a non-environmentally friendly issue. But we should require our customer service department to set up a complaint feedback mechanism to deal with the potential cases of this issue. For example, if a customer realises that the product is harmful because of the materials and requests it to be returned, we should respond quickly with sincere apologies and issue a refund.

\subsection{Make Long-Term Commitments to Sustainability}

We made a mistake in selecting materials, and to make amends, we must commit to long-term sustainability in our manufacturing processes. Sustainability is the foundation of our industries and without it there would be no more marketing or manufacturing. We must continue to seek raw materials which are really environmentally friendly and recyclable. We should keep our promise in front of our consumers.

\subsection{Establish Marketing Oversights}

We should set up a work group of marketing oversight under the marketing department to avoid issues like false advertising and misleading. It is claimed that advertising ethics are important to a company's long-term development \citep{schauster2018ethics}. These oversight people should not be involved in advertising decisions but just give their opinions on the outcomes by their colleagues. Once the oversight team believes that there are false statements in our advertising, we must remake our decisions in the earlier advertisement.

\subsection{Improve the Raw Materials Examination}

After making the commitments, we should improve the examination of raw materials because we were using non-environmentally friendly materials, which proves that there must be a defect in our materials choice. A good way to recognise the materials is third-party testing agency engagement. If we could not examine the materials by ourselves, we could let the professionals do it for us, which could reduce the harm of using non-environmentally friendly materials.

\subsection{Conduct Regular Environmental Impact Assessments}

Besides, we could also conduct an assessment of our product's environmental impact. Making sure not only the manufacturing processes are environmentally friendly is important, but also the outcomes should not do harm to the environment. We should also engage a third-party agency to assess our products, ensuring transparency and demonstrating our commitment to never repeating the same mistake. If the assessment shows that our product is still not meeting the environmental standards, we should suspend production for rectification.

\subsection{Enhance the Ethical Guidelines}

After this case, what we learn is to avoid using non-environmentally friendly materials in production. We could transform the case consequences into valuable lessons and add them into our ethical guidelines. We must place a strong emphasis on using environmentally friendly materials in future employees' training to avoid the same mistake.

\texttt{Words: 1650}

\section{Summary and Conclusions}

Finally, let us organise the outline of the whole case. We are a manufacturer of a product with raw materials which are not environmentally friendly. But what we claimed is that we are totally using environmentally friendly materials in the process. These behaviours caused ethical violations which are losing public trust and lacking transparency. These violations affect about 15,000 people and do real damage to the ocean environment and may cause overexploitation of resources. Besides, this issue does cause a lot of ethical influences to our company including reputation reduction, fines and compliance risks.

As the crisis rises, what we can do is hold transparent and clear communication with our consumers to explain the cause and effect. Next we should initiate the recall procedure and set up a complaint and feedback system. Then we should make a commitment to using environmentally friendly materials in our manufacturing, with real actions like establishing oversight teams and conducting regular environmental impact assessments.

What we have learned is that humankind's development should go with environmental sustainability. If we leave the environment alone and let us exploit the resources, there will soon be nothing for us and our future generations to use. Another important lesson is that we must show honesty in our marketing, and we should advertise our products from facts but not false statements. Once we pursue perfection in these fields, our products could become popular among customers.

\bibliographystyle{agsm}
\bibliography{report_references}

\appendix
\section{Appendices}

\subsection{Pitch Video}

\href{https://youtube.com/shorts/PAeucE88hLk}{Pitch Link}

\subsection{Pitch Slides}

\includepdf[
  pages=-,
  scale=0.9,
  pagecommand={\thispagestyle{plain}}
]{PitchSlides.pdf}

\end{document}
